\begingroup
\fontsize{8.5pt}{9.5pt}\selectfont
\renewcommand{\arraystretch}{1.3}
\setlength{\tabcolsep}{4pt}

\begin{longtable}{|
	>{\raggedright\arraybackslash}p{4.0cm}|
	>{\raggedright\arraybackslash}p{9.5cm}|}

	\caption{Estructura de los casos de prueba}
	\label{tab:estructura-casos-prueba}
	\\
	\hline
	\tableheader
	\textbf{Campo}             &
	\textbf{Descripción}
	\\
	\hline
	\endfirsthead

	\hline
	\tableheader
	\textbf{Campo}             &
	\textbf{Descripción}
	\\
	\hline
	\endhead

	ID                         & Identificador único asignado al caso de prueba.                                                                \\ \hline
	Nombre                     & Denominación que describe de forma breve el comportamiento evaluado.                                           \\ \hline
	Requisitos asociados       & Requisitos que pueden ser evaluados mediante la prueba.                                                        \\ \hline
	Objetivo                   & Comportamiento específico que se pretende verificar.                                                           \\ \hline
	Resultado esperado         & Comportamiento que debería producirse si la solución funciona de acuerdo con lo establecido.                   \\ \hline
	Precondiciones             & Condiciones que deben cumplirse antes de iniciar la prueba.                                                    \\ \hline
	Datos de prueba            & Usuarios, grupos, credenciales u otros elementos utilizados durante la ejecución.                              \\ \hline
	Procedimiento y evidencias & Secuencia de los procedimientos realizados, indicando cada literal y acompañándolo de su respectiva evidencia. \\ \hline
	Estado                     & Clasificación final de la prueba como cumple, cumple parcialmente o no cumple.                                 \\ \hline
	Observaciones              & Información adicional relevante sobre la ejecución, incidencias o limitaciones identificadas.                  \\ \hline
\end{longtable}
\endgroup
\normalsize
